\documentclass[aps,prd,twocolumn,superscriptaddress,nofootinbib]{revtex4-2}

\usepackage{amsmath,amssymb,amsthm}
\usepackage{graphicx}
\usepackage{bm}
\usepackage{booktabs}
\usepackage{xcolor}
\usepackage{hyperref}
\usepackage{orcidlink}
\usepackage{mdframed}

\newtheorem{axiom}{Axiom}
\newtheorem{theorem}{Theorem}
\newtheorem{proposition}{Proposition}
\newtheorem{corollary}{Corollary}
\newtheorem{definition}{Definition}
\newtheorem{remark}{Remark}

\newcommand{\lean}[1]{\texttt{\small #1}}
\newcommand{\SM}{Supplemental Material}

\begin{document}

\title{Scale-Coupled Dynamics: Gravity and Charge from a Single Non-Commuting Scale Field}

\author{Xinyu Yang\,\orcidlink{0009-0007-2600-0948}}
\noaffiliation

\date{\today}

\begin{abstract}
We develop an axiomatic theory in which \emph{scale} — a dimensionless
local unit ratio — rather than spacetime is primitive. From seven axioms
we derive that the local scale must be directional (a scalar version is
refuted by its own consequences), that two scalings along different
directions must fail to commute with an antisymmetric residue that is
exactly a rotation, that the resulting Lie algebra has real rank one and
hence belongs to the Lorentz family, and that the number of directions is
fixed by a single postulate about what counts as an observation.
Curvature, the gauge connection, and the renormalization group turn out
to be different readings of one object — the scale vector field and its
own rate of change. The isotropic sector reproduces the Schwarzschild
metric and closes quantitatively: light bending $1.7515''$ and Mercury's
perihelion advance $42.989''/\mathrm{century}$, both to the precision
quoted, with no free parameter beyond the single dimensionless unit any
dimensionless theory must carry. The anisotropic sector produces point
particles as unresolvable knots of the scale vector field, carrying
exactly a $\mathbb{Z}$ winding charge and a $\mathbb{Z}/2$ confined
charge. Non-commutativity of the scale generators produces $\hbar$ as an
exact step along the scale axis and makes the entire quantum correction
to curvature purely antisymmetric, proving that the radiative sector of
gravitational waves is identical to general relativity at \emph{every}
post-Newtonian order — the two theories differ, if at all, only through
the absence of a horizon. We confront the framework with GW250114's
area-law test, the ergoregion-instability constraint from rapidly
spinning remnants, DESI's dark-energy equation of state, the $f_0(500)$
pole width, and fourteen decades of the measured particle spectrum, and
report honestly where the framework currently loses: a
$2.8$–$4.2\sigma$ tension with DESI's preference for an evolving equation
of state, and an open question at the $f_0(500)$ pole. The argument has
been verified end to end in the Lean~4 proof assistant — 644
machine-checked theorems, zero admitted gaps — reproduced as ordinary
mathematics in the \SM, which this self-contained paper does not require
to be read.
\end{abstract}

\maketitle

%======================================================================
\section{Introduction}\label{sec:intro}
%======================================================================

General relativity and the Standard Model both take a four-dimensional
spacetime manifold as given and place fields upon it. Each has open
seams at the place where it meets the other, or meets its own
foundations: a singularity where curvature diverges and the classical
description simply stops; a cosmological constant whose value is
measured, not explained, by either theory; a renormalization group
whose running couplings are computed order by order in perturbation
theory rather than obtained in closed form; a particle spectrum whose
labels — charge, color, generation — are read off experiment and then
built into the group-theoretic structure of the Lagrangian by hand. None
of these are failures of either theory on its own terms; they are the
places where a choice was made — a manifold, a gauge group, a set of
fields — that the theory itself does not derive.

This paper explores a different starting point, one deliberately more
austere than either: take the \emph{scale} at a point — a dimensionless
ratio of a local unit to a reference — as the only primitive, and ask
what follows if that scale is directional, because the physical sources
that produce it (energy-momentum) are tensorial rather than scalar, and
non-commuting, because performing two scalings along different
directions in the wrong order leaves a residue. Spacetime curvature, the
gauge connection, particle content, and $\hbar$ are then not
independently postulated; they are different questions asked of a single
object, the scale vector field, and several of the open seams listed
above turn out to be forced closed by that object's own algebra rather
than left as free choices: the dimension and signature of spacetime, the
particle-label alphabet, the leading term of the renormalization group,
and the relation between the metric connection and the gauge connection
all become theorems rather than inputs.

This is not a claim that spacetime does not exist, nor a claim that the
program competes with approaches that quantize the metric directly (loop
quantum gravity), posit extended objects and extra dimensions (string
theory), or build spacetime from discrete causal order (causal set
theory). Those programs start by choosing what is fundamental among
geometry, matter, and causal structure, and then quantize or discretize
it. This paper asks a prior question: whether geometry, gauge structure,
and particle content are themselves derivable from something more
primitive than any of the three — and answers it constructively, by
exhibiting the derivation rather than arguing for its possibility in the
abstract. The claim is that the four-dimensional, Lorentzian, curved
structure usually assumed can instead be \emph{derived} from a single
forced algebraic fact: two symmetric operators bracket to an
antisymmetric one. Everything that follows is the working-out of that
one fact's consequences, and nothing else is smuggled in along the way —
a claim that Section~\ref{sec:register} exists specifically to make
checkable rather than asserted.

Every derivation summarized here has an independent, complete proof; we
have verified the entire argument mechanically in the Lean~4 proof
assistant against no axioms beyond those of standard mathematics
(\texttt{propext}, \texttt{Classical.choice}, and \texttt{Quot.sound}),
and the \SM\ reproduces every statement and proof — 644 theorems in
total, spanning definitions, lemmas, and closing numerical bounds — in
full mathematical detail, organized to mirror this paper section by
section. This paper does not depend on that verification to be read; it
is a self-contained physics argument, written so that a reader who
never opens the \SM\ still receives every derivation's actual content
rather than a pointer to it. What the verification adds is a guarantee
that cannot be obtained by proofreading alone: that every ``Theorem'' in
what follows is exactly what it claims to be, with every hypothesis
stated and no step skipped, because a proof assistant does not accept an
argument it cannot check line by line.

The remainder of the paper is organized as follows. Section~\ref{sec:picture}
states the physical picture in ordinary language before any algebra,
because a statement built from commutators and a Lie-algebra
classification is not obviously less abstract than the curved geometry
it replaces, unless a picture is attached to it first; the picture
introduces, informally, every structural result the later sections
derive formally, so that the algebra which follows always has a physical
question attached to it. Section~\ref{sec:axioms} states the seven
axioms precisely. Section~\ref{sec:chain} derives the forced chain from
the axioms to a four-dimensional Lorentzian structure carrying both a
metric connection and a gauge connection — the single argument on which
the rest of the paper depends. Section~\ref{sec:branch} states the
branch point between the isotropic and anisotropic sectors, which
determines whether a given region of the scale field carries curvature
alone or curvature together with charge. Section~\ref{sec:gravity} gives
the isotropic (gravitational) sector in full: the Schwarzschild solution
and the two classical solar-system tests, the Newtonian limit obtained
exactly rather than as a truncated series, local covariance under global
dissipation, the absence of singularities and horizons, and
gravitational-wave dynamics from inspiral through ringdown, including a
concrete confrontation with GW250114~\cite{GW250114} and
GW241011~\cite{GW241011}.
Section~\ref{sec:particles} gives the anisotropic (particle) sector: the
two topological charges, their independence, and the resulting bound on
the mass spectrum. Section~\ref{sec:flow} gives the renormalization
group, obtained as a theorem about a defect density rather than a
perturbative computation, and the strong-interaction scale.
Section~\ref{sec:quantum} gives the quantum sector — $\hbar$ as an exact
step along the scale axis — and the all-orders equivalence with general
relativity in the gravitational-wave sector that follows from it.
Section~\ref{sec:cosmology} gives dark matter, dark energy, and the
scale-blindness of light. Section~\ref{sec:data} confronts the framework
with data, in one table. Section~\ref{sec:falsify} states, observation by
observation, what would kill which part of the framework.
Section~\ref{sec:register} is a register of what has been retracted,
corrected, cited without proof, and assumed — the standing answer to
whether any of the foregoing rests on more than the seven axioms plus one
external physical input. Section~\ref{sec:concl} concludes.

%======================================================================
\section{The physical picture}\label{sec:picture}
%======================================================================

Sixteen statements, in the order they compose, with the algebra deferred
to the sections that follow. Each previews a result proved formally
later; the section reference at the end of a paragraph is where the
informal claim becomes a theorem.

\emph{There is no space, only a ratio.} At every point there is a
number, and that number is a ratio: the local unit compared to some
reference. The choice of reference is unobservable by construction, so
only \emph{differences} of this number are meaningful. Picture a map
that records, everywhere, how long a meter is there — with no grid drawn
on it. The grid, i.e.\ spacetime, is what emerges from the
\emph{inconsistency} between these numbers at nearby points, not
something the numbers are written onto. This is a stronger claim than
``spacetime is emergent'' in the sense usually meant by that phrase: it
is not that a smooth manifold approximates some discrete substrate at
long wavelengths, but that the manifold structure itself — the fact that
nearby points have well-defined relations at all — is read off
differences in a single scalar-valued field, with no independent
metrical input.

\emph{``Two places'' is constructed, not primitive.} If the scale
pattern at two points is parallel, no observable distinguishes them —
not two places that look alike, but no second place at all, since every
question one could ask of one point has the identical answer at the
other. A perfectly uniform region has no internal structure and no
extent: distance is the accumulated \emph{non-parallelism} between
points, and because the wedge product of two vectors is antisymmetric,
separation is naturally a signed area, not a length. Two consequences
follow immediately, both proved in Section~\ref{sec:chain}: a truly
homogeneous universe would be a single point no matter how large a
coordinate volume one assigns it, and the ordinary notion of distance as
a positive scalar is recovered only after choosing an orientation, which
this framework does not supply for free.

\emph{Scale is directional because energy-momentum is tensorial.} At
each point the scale is not a number but a \emph{vector}: a fluid
carries different energy density along its flow than transverse to it,
and a tensorial source cannot be matched by a single scalar scale. This
is not an aesthetic preference for more structure; Section~\ref{sec:gravity}
proves the scalar version of this requirement is not too weak but too
strong — it \emph{forbids} the Schwarzschild metric outright, so a
scalar scale is not merely insufficient to describe the solar system, it
actively excludes the possibility of describing it. Directionality is
consequently not an enrichment chosen for generality; it is the minimum
needed for the theory to be compatible with a single measured orbit.

\emph{Gravity is the ruler changing, not space bending.} A change in the
\emph{magnitude} of the scale vector across position is exactly gravity;
curvature is the second-order inhomogeneity of scale — not a first-order
gradient, which is unobservable, but the rate at which that gradient
itself changes. Nothing is bent: light follows what looks locally
straight, but the ruler used to judge straightness changes along the
path, so the trajectory appears curved. Tidal forces are the
\emph{second}-order change in the ruler — the first-order, uniform-dilation
part is unobservable, being exactly a reference shift. Two bodies in
free fall separate or approach not because a force acts between them but
because the local unit of length differs infinitesimally at their two
positions, and that difference is what a tidal-force meter actually
measures; the entire content of the gravitational field, on this
reading, is a statement about how a ruler's length varies from place to
place, encoded completely in the second derivative of one scalar field
per direction (Section~\ref{sec:gravity}).

\emph{Rotation is direction changing — and it is \emph{generated} by
scaling.} A change in the \emph{direction} of the scale vector across
position is rotation, and here is the theory's sharpest forced step:
performing two scalings along different directions in the two possible
orders must differ by a rotation. This is not a chosen coupling; it is
forced by the meaning of the words — a scaling is a pure stretch and
hence symmetric, a rotation preserves length and hence is antisymmetric,
and the bracket of two symmetric operators is necessarily antisymmetric
(Theorem~\ref{thm:coupling}). Physically this is the Thomas–Wigner
rotation~\cite{Thomas1926}, and the
Terrell–Penrose effect~\cite{Terrell1959,Penrose1959} is its visual form: a
fast-moving object is \emph{seen} as rotated rather than contracted,
because a change of scale is being read as a rotation. One direction:
flat. Two or more: curvature, generated by the scalings themselves, with
no separate assumption about how gravity and rotation relate — they are
the symmetric and antisymmetric parts of one algebra, and the bracket
between them is arithmetic, not physics chosen to fit.

\emph{The gauge connection and the gravitational connection are the same
object.} Both are built from the same log-derivative of the scale
pattern, evaluated on two different loci of the same field: where the
pattern is directionally uniform (isotropic), the connection reads as
curvature; where it is not (anisotropic), the same construction reads as
a Yang–Mills field strength. A gauge theory and a metric theory are, in
this framework, two questions asked of one connection, not two
structures glued together after the fact (Section~\ref{sec:chain}).

\emph{``Unifying the forces'' is a question in a vocabulary this theory
does not have.} The standard unification program asks for one gauge
group from which every force descends, which presupposes that forces
are couplings and that couplings are the fundamental description; in
this framework neither holds, since comparisons of local units are
primitive and forces are read off them. The question this framework can
answer instead is what is expressible at all: a quantity is observable —
invariant under A5's fiducial shift — exactly when it is a function of
scale \emph{differences}, not merely typically or approximately, so the
scale channel expresses differences and nothing else, and an absolute
scale value is consequently not expressible by any construction
available here, not merely unmeasured. This is also the reason no
single cross-sector test between the scale channel and the rotation
channel exists: in the fully scalar reading the two channels are
independent by construction, one being an abelianization and the other
its kernel, so neither can constrain the other — and Section~\ref{sec:chain}'s
directional coupling is exactly what breaks that independence once
scale is given a direction, replacing "independent channels" with "one
channel generating the other."

\emph{Whether charge exists at all is itself a fact about scale.} The
branch between the isotropic and anisotropic loci is not fixed once and
for all; it is a property of the scale pattern at a given resolution,
and the same region of the substrate can sit on one side of the branch
at one energy scale and the other side at a different one. This is the
framework's reading of the empirical fact that interactions separate as
energy falls: it is not that several independent forces happen to
decouple at different scales, but that whether an order parameter, and
hence a charge, is present at all is a scale-dependent question with a
single underlying source (Section~\ref{sec:branch}).

\emph{$\hbar$ is a step along the scale axis.} Not a mysterious
constant: translating one step along the scale axis fails to commute
with position, and the size of the mismatch is $\hbar$ — exactly, not as
a first-order approximation in some formal deformation parameter that is
later set to a physical value. Uncertainty is then not measurement
disturbance but the mutual exclusion of scale sharpness and content
richness: a perfectly sharp scale is a single point on the scale axis,
and a point carries no content, so resolving any finite amount of
structure requires a finite window, and the width of that window times
the density of what it must resolve is bounded below
(Section~\ref{sec:quantum}).

\emph{The entire quantum correction to geometry is a rotation, never a
stretch.} Once the scale algebra is made non-commutative, the correction
it produces to curvature is exactly the commutator of two scale
generators — and a commutator, by the same symmetric/antisymmetric split
that produces gravity and rotation in the classical theory, is purely
antisymmetric. The symmetric, radiative part of the geometry is
consequently untouched by quantization to any order: gravitational waves
are the one observable sector in which this framework and general
relativity cannot differ, not merely one in which the difference is
currently too small to see (Section~\ref{sec:quantum}).

\emph{``What exists'' is a slice, not a property of the world.} Changing
resolution does not mean knowing more about the same particles; it means
looking at a different set of them. ``The list of fundamental
particles'' is accordingly ill-posed: there is always unresolved content
above any scale, with no final layer. The size of the Standard Model's
label set measures the logarithmic range it covers, not physical depth
— a theory that must remain valid over thirty orders of magnitude in
energy needs a label for everything that appears anywhere in that range,
whether or not it is realized at any particular scale one happens to
probe.

\emph{A particle is an unresolvable knot in the scale vector field.} Not
a field excitation but a configuration: the vector winds around a point
and cannot be combed flat. Its integer charge is the winding number of
the vector around a sphere; its mass is not a property it carries but
the location at which it appears on the scale axis; its rotational label
is $v \wedge \nabla v$, so charge necessarily comes with spin. A second,
independent label — a confinement class rather than a winding number —
is available from the same construction, and the two are logically
independent of each other: knowing one fixes nothing about the other
(Section~\ref{sec:particles}). Exactly two labels are available, not
three and not one, because the scale pattern in the anisotropic sector
has already been shown to collapse to a single vector
(Section~\ref{sec:chain}), and a single vector supports exactly a
winding number and a monodromy class and nothing further.

\emph{The renormalization group is a theorem about counting, not a
perturbative expansion.} A coupling read as a ratio of defect densities
changes with the logarithm of scale at a rate fixed directly by the same
dissipation constant that drives cosmic expansion (Section~\ref{sec:flow}).
No diagram is summed and no loop order is chosen; the leading term of
the beta function is an algebraic consequence of what a coupling
\emph{is} in this framework; a subject usually organized around
successive orders of a perturbative expansion is organized here around a
single exact linear relation.

\emph{The universe is not expanding; scale is running.} Galaxies are not
flying apart — the ruler is shrinking, and because energy scale is the
reciprocal of length scale, the energy scale is dissipating. There is no
first moment, because there is no zero scale: scale is a unit and is
never zero, so ``the age of the universe'' is not a quantity — what
exists is an elapsed ratio, finite and measurable, with no origin
(Section~\ref{sec:cosmology}).

\emph{Black holes are arbitrarily dark, never fully black.} The horizon
is where the time component of the metric vanishes, and that component
is minus the square of a scale unit — never zero. No singularity, no
horizon, no infinite redshift, all three the same fact: ratios are
invertible. A black hole is a region of extreme but finite scale ratio;
the exterior geometry is untouched, so shadow, ISCO, and the early
ringdown all agree with general relativity, but nothing is causally cut
off and no information is destroyed (Section~\ref{sec:gravity}).

\emph{A horizonless remnant should ring the way a black hole does, then
echo.} Because nothing in the metric is ever exactly reflecting or
exactly absorbing, a gravitational-wave signal that would terminate at a
horizon in general relativity instead partially reflects back and forth
in this framework, producing a train of echoes evenly spaced in time
after the main ringdown, with amplitude falling geometrically from one
echo to the next. This is the one place this framework's predictions are
qualitatively distinctive without being quantitatively computable, since
the reflectivity itself depends on physics — the interior equation of
state — that is explicitly outside this framework's scope
(Section~\ref{sec:gravity}).

\emph{The two arrows of time are one order.} A long-standing puzzle
elsewhere is why the direction of entropy increase coincides with the
direction of cosmic expansion; the usual answer invokes a special
initial condition aligning two otherwise independent facts. Here they
are not two facts. Three results, proved independently and never
invoking one another — that time is the direction along which scale
drifts (Section~\ref{sec:chain}), that an $H$-theorem for coarse-grained
content holds using only finite resolution, with no scale dependence
anywhere in its own proof, and that the entropy of a region of scale
ratio $R$ is $\rho \ln R$ (Section~\ref{sec:cosmology}) — combine into
the following.

\begin{theorem}[The two arrows are one order]\label{thm:onearrow}
For $\rho>0$ and $R_1,R_2>0$,
\[
\rho\ln R_1 \le \rho\ln R_2 \iff R_1\le R_2,
\]
and likewise with strict inequalities.
\end{theorem}
\begin{proof}
$\rho \ln$ is a strictly increasing function of $R$ for fixed $\rho > 0$,
since $\ln$ is strictly increasing and multiplication by a positive
constant preserves order.
\end{proof}

Ranking regions by entropy and ranking them by scale ratio are
\emph{literally} the same order relation, not two arrows that happen to
align. The one modeling choice this rests on is reading entropy as
$\rho \ln R$, the count of resolvable thresholds within a region of scale
ratio $R$; that reading is motivated by, but not forced by, the
threshold-density construction of Section~\ref{sec:flow}, so the identity
is exact given the reading and the reading itself is a choice, stated as
such.

\emph{Summary.} A number (ratio) becomes a vector because its source is
tensorial; the vector's magnitude changing is gravity, its direction
changing is rotation; the two are coupled by the meaning of the words
involved; the commutator of that coupling gives $\hbar$ and the quantum
correction; where the vector cannot be combed flat is a particle; the
vector running overall is cosmology. All physics here is a different
reading of the same vector field. What this picture does \emph{not} buy
is stated equally plainly: no coupling constants, no mass values, no
generation structure, no scattering amplitudes. It answers ``what,'' not
``how much''; the quantitative part closes fully only in the
gravitational sector, and the particle sector delivers label structure
and bounds, not numbers.

%======================================================================
\section{Axioms}\label{sec:axioms}
%======================================================================

The substrate is a commutative differential ring over $n$ directions; an
earlier version of this development posited three distinct carrier
structures for different sectors of the theory, later shown to reduce to
one. Seven axioms are used throughout.

\begin{axiom}[A1 — substrate]
There is a commutative ring $A$ with $n$ commuting derivations
$d_1,\dots,d_n : A \to A$, each additive and satisfying the Leibniz rule.
\end{axiom}

\begin{axiom}[A2 — scale field]
A distinguished unit $s \in A^\times$, the local scale.
\end{axiom}

\begin{axiom}[A3 — scale-energy duality]
An energy unit $\varepsilon\in A^\times$ with $\varepsilon\cdot s = 1$.
\end{axiom}

\begin{axiom}[A4$'$ — directional measurement]
Scale is assigned per direction: $s : \{1,\dots,n\}\to A^\times$, not a
single scalar. (An earlier, scalar version, A4, is refuted in
Section~\ref{sec:gravity}: it forbids the Schwarzschild metric outright.)
\end{axiom}

\begin{axiom}[A5 — fiducial covariance]
Physical content is invariant under $s \mapsto u\cdot s$ for any global
unit $u$: only ratios of scale are observable.
\end{axiom}

\begin{axiom}[A6 — openness / dissipation]
The scale field is not conserved by a global potential: no functional's
gradient equals the full drift.
\end{axiom}

\begin{axiom}[A7 — observability]
A commitment about what counts as a single observation, formalized as a
dimension-matching condition between the rotation sector the coupling
generates and the direction sector A4$'$ labels; it fixes the number of
\emph{independent} directions rather than leaving it a free integer.
\end{axiom}

A7 is deliberately an axiom and not an informal judgment: it is a
commitment about what an observation is rather than a claim about the
world, and Section~\ref{sec:falsify} returns to the consequence that it
cannot be falsified by any single experiment.

One further ingredient, the scale response law
\begin{equation}
e^{2\sigma}R = \kappa \rho, \tag{A6$'$}
\end{equation}
relating the curvature scalar $R$ built from the scale field to a matter
density $\rho$, is used in the gravity sector. \textbf{It is not one of
the seven axioms.} It is an external physical input — the one place
general relativity's own field equation is assumed rather than derived —
and it is registered as such in Section~\ref{sec:register}.

%======================================================================
\section{The forced chain to a four-dimensional Lorentzian gauge structure}\label{sec:chain}
%======================================================================

This section is the spine of the paper: a single chain of implications,
each step forced rather than chosen, from A1–A7 to a four-dimensional
Lorentzian structure carrying both a metric and a gauge connection.

\textbf{The scale/rotation split is not a choice.} Given A1–A4$'$, any
endomorphism of the tangent data at a point decomposes canonically into a
symmetric part (a scaling) and an antisymmetric part (a rotation); this
is a fact about the algebra of endomorphisms of an inner-product space,
not a modeling decision, since every such endomorphism is the sum of its
symmetric and antisymmetric parts and that decomposition is unique. An
earlier stage of this development reached the same split by a different
route, an abelianization of an abstract comparison group applicable only
to the scalar case; that construction is superseded by, and is a special
case of, the endomorphism decomposition used from here on, and the two
are not independently combined into a stronger joint claim.

\textbf{Time is the direction along which scale drifts.} Given A6
(dissipation), exactly one direction carries a definite, non-cancelling
sign of drift; the signature is forced to a codimension-one split,
$(1,n-1)$. A theory in which every direction could drift in either sign
with no distinguished one would admit no consistent notion of ``before''
and ``after''; A6 rules that out. Labeling the distinguished direction
\emph{timelike} rather than merely \emph{distinguished} is a convention
consistent with the split, not a further consequence A6 forces on its
own — a distinction registered explicitly rather than smoothed over
(Section~\ref{sec:register}).

\textbf{The bookkeeping form is $\tau(xy)$ for a trace $\tau$ alone,} not
an independent postulate: any bilinear form on the ring compatible with
both the ring product and the derivations must take this shape, because a
trace is (up to scale) the unique linear functional invariant under
conjugation, and conjugation invariance is what compatibility with the
derivations demands.

\textbf{Curvature is the commutator of transports.} Let $\nabla_i$ denote
parallel transport of a tensorial quantity along direction $i$, built
directly from the scale pattern rather than postulated independently.
Then the curvature of the connection is
\begin{equation}
F_{ij} := [\nabla_i, \nabla_j],
\end{equation}
literally the failure of two transports to commute, rather than an
independently defined tensor whose vanishing is imposed or assumed. Two
consequences follow directly from this definition rather than from extra
input: first, commuting \emph{derivations} ($d_id_j = d_jd_i$, part of
A1) does not imply a flat \emph{connection} — $F_{ij}$ can be nonzero
even when the underlying derivations commute, because curvature lives in
the connection, not in the chart; second, the same commutator $F_{ij}$ is
what is computed in the gravitational sector as the Riemann curvature and
what is computed in the anisotropic sector as the Yang–Mills field
strength (Section~\ref{sec:particles}) — \textbf{gravity and gauge
structure are the same object, read under two different questions about
which sector's scale pattern is being transported.} Neither is assumed;
both are the same commutator. This identity extends one step further,
to the quantum sector of Section~\ref{sec:quantum}: the commutator
$[D_i,D_j]$ that gives curvature here is the identical construction
that gives $\hbar$ there, evaluated on a different family of
transports — phase-space translations in one case, frame transports in
the other. Curvature and $\hbar$ are consequently not merely analogous;
they are one algebraic construction read on two different transport
families, a stronger and more specific claim than any analogy between
gravity and quantum mechanics usually offered.

A further split, orthogonal to the isotropic/anisotropic branch of
Section~\ref{sec:branch}, separates transport itself into a scale part
and a frame-rotation part. The scale part, being a gradient by A5, is
automatically integrable — mixed partial derivatives commute — and
contributes nothing whatsoever to the curvature $F_{ij}$; all of it is
carried by the frame-rotation part. This is the structural reason this
framework avoids the defect that killed Weyl's original 1918 unified
theory~\cite{Weyl1918}, where an analogous non-Riemannian connection
produced a length that depended on the path taken between two points (a
``second clock effect'' ruled out by observation): here the scale
transports integrably by construction, so no such path-dependence can
arise, while curvature — gravity — lives entirely in the
non-commuting frame-rotation sector. Gravity is consequently not, even
informally, ``the scale field bending something''; it is the failure of
frame rotations specifically to commute, with the scale itself never
contributing a discrepancy of its own.

\textbf{Directions are the defining representation, not the algebra
itself.} Once the algebra generated by the scalings and rotations is
named explicitly, its dimension as a Lie algebra ($\dim\mathrm{so}(3,1)=6$
for the eventual case $n=4$) exceeds the number of directions ($4$), so
directions form the vector representation on which the algebra acts, not
the algebra's own coordinates — a distinction an earlier version of this
development conflated.

\begin{theorem}[Coupling is forced]\label{thm:coupling}
Let $\Sigma_i$ be the scaling generator along direction $i$. Then
$[\Sigma_i,\Sigma_j]$ is nonzero for $i\neq j$ and is purely
antisymmetric, i.e.\ a rotation generator.
\end{theorem}
\begin{proof}
A scaling is symmetric with respect to the pairing induced by the trace
form (it preserves direction and only rescales magnitude, so it is
self-adjoint); a rotation is antisymmetric (it preserves length, so its
generator satisfies $R^T = -R$). For symmetric $A,B$,
\[
(AB-BA)^T = B^TA^T - A^TB^T = BA - AB = -(AB-BA),
\]
so the commutator of two symmetric operators is antisymmetric. That the
bracket is additionally \emph{nonzero} for $i\neq j$ — that the coupling
is not merely well-defined but genuinely active — is witnessed concretely
rather than argued in general: an explicit boost pair in three spatial
dimensions exhibits a nonvanishing bracket, and no general
nondegeneracy argument beyond that concrete witness is claimed.
\end{proof}

No coupling constant is chosen here; the coupling exists purely because
of what the words ``scaling'' and ``rotation'' mean algebraically. The
generated rotation is not merely known to exist and be antisymmetric;
its concrete form is fixed once the two scaling directions are: for
scalings along direction vectors $v,w$, the rotation their bracket
generates is exactly the antisymmetric bilinear pairing
$(v\wedge w)_{ij} = v_iw_j-v_jw_i$, the wedge of the two directions.
This makes the coupling's output a genuine two-form — antisymmetric,
bilinear, and vanishing exactly when the two scaling directions are
parallel — rather than an abstract nonzero element asserted only to
exist; a wound configuration in the anisotropic sector, whose direction
genuinely varies rather than staying parallel to itself, consequently
carries this rotational label automatically, while a homogeneous
pattern — one direction wedged with itself — carries none, whatever
magnitude it has.

\textbf{Real rank one, hence the Lorentz family.} The Lie algebra
generated by one vector's worth of scalings together with the rotations
they force has real rank one: the maximal abelian subalgebra of
diagonalizable elements is one-dimensional, since a second independent
scaling direction would, by Theorem~\ref{thm:coupling}, generate a
rotation mixing it with the first, destroying simultaneous
diagonalizability. Citing, without reproving, the classification of
real-rank-one real simple Lie algebras — $\mathrm{so}(n,1)$,
$\mathrm{su}(n,1)$, $\mathrm{sp}(n,1)$, $\mathrm{f}_{4(-20)}$ — this
places the algebra in the Lorentz family $\mathrm{so}(n,1)$; the other
branches require additional structure (a complex, quaternionic, or
octonionic module structure on the directions) absent from a real
scale-vector construction, leaving $\mathrm{so}(n,1)$ as the surviving
case. The step from ``one drift functional'' to ``one-dimensional
maximal abelian subalgebra'' above is a physical identification of what
counts as a single independent scaling direction, not a formal
deduction from A1–A7 alone — stated here rather than left implicit, since
it is exactly the kind of step where a numerical coincidence can hide a
choice (Section~\ref{sec:register}).

\textbf{$n=k+1$, and $k=3$ from A7.} The number of directions $n$ is one
(time) plus the number of independent spatial directions $k$. $k$ is not
free once A7's dimension-matching condition is imposed: requiring the
number of rotations the coupling generates to match the number of
directions A4$'$ labels gives $k=3$. This is not the only matching
condition one could write down — comparing the dimension of the full
algebra to the dimension of its own defining representation instead
gives $k=2$, and the two conditions provably disagree — so A7 is doing
real work: it is a specific commitment about \emph{which} two quantities
must be equal, not an arithmetic inevitability. The framework's own
argument for preferring the first condition is that the labels available
are the directions A4$'$ actually assigns scales to, and the coupling's
output must match \emph{those}, not the abstract size of the algebra
that generates it — but this is an argument for a modeling choice, not a
proof that no other choice is coherent, and it is why A7 is registered
as a postulate about what counts as an observation rather than presented
as forced by pure counting.

\textbf{The pattern is one vector; more axes are unavailable.} An
apparent dilemma between a ``uniaxial'' and a ``biaxial'' reading of an
earlier theorem is resolved once the relevant vanishing identity,
$v \wedge v = 0$, is recognized as a statement about \emph{homogeneity} —
a vector wedged with itself is always zero, regardless of how many axes
exist — rather than a constraint on the number of axes. Once that
misreading is removed, the scale pattern is shown to collapse to exactly
one vector, and additional independent axes are not merely disfavored
but structurally unavailable given the coupling of Theorem~\ref{thm:coupling}.

The chain therefore delivers, from A1–A7 with no further input: a
Lorentzian signature, four dimensions, a metric and a gauge connection
that are the same underlying commutator, and a coupling between scaling
and rotation forced by algebra rather than fit to data.

%======================================================================
\section{The branch point}\label{sec:branch}
%======================================================================

The scale pattern at a point is either \emph{isotropic} — the same in
every direction — or \emph{anisotropic}, and whether it is one or the
other is a scale-dependent fact, not a fixed property of the point:

\begin{center}
\begin{tabular}{ll}
\toprule
isotropic & anisotropic \\
\midrule
no order parameter & order parameter present \\
no charge & charges present \\
no multiplet & multiplets present \\
no rotation-sector charge & curvature/rotation active \\
\bottomrule
\end{tabular}
\end{center}

Both columns have one source: isotropy is a single algebraic condition
on the pattern (uniform direction across the pattern's components), and
charge, multiplet structure, and rotation-sector activity all switch on
or off together with it, not independently. The formal branch condition
itself is a static property of a given pattern, not a function of a
separate energy-scale variable; the further, motivating reading — that
the same physical region can present as isotropic at one probing
resolution and anisotropic at another, giving this framework's account
of interactions separating as energy falls — is an interpretation
connecting this branch to the scale-flow structure of
Section~\ref{sec:flow}, not itself a theorem proved from the branch
condition alone, and is stated here as such. Section~\ref{sec:gravity}
treats the isotropic sector; Section~\ref{sec:particles} the anisotropic
one.

An apparent contradiction is worth forestalling: Section~\ref{sec:gravity}
computes curvature in a sector called ``isotropic.'' ``Isotropic'' here
means uniform \emph{direction} of the scale vector, not uniform
\emph{magnitude} — a radially varying but directionally uniform scale
pattern has no charge and no rotation-sector label, and still has
curvature, because curvature is generated by magnitude inhomogeneity
(Section~\ref{sec:picture}), which is an axis independent of direction
uniformity.

A second consequence of the branch structure fixes what happens
\emph{inside} a multiplet, \emph{given} an assignment of one scaling
operator to each direction determined by that direction's own scale
value — a hypothesis satisfied automatically whenever such an
assignment exists, though real rank one (Section~\ref{sec:chain}) only
guarantees one vector and one scaling overall, not a free per-direction
assignment, so the statement below is conditional rather than a further
automatic consequence of the chain. Under that hypothesis, two
directions assigned to the same multiplet cannot rotate into each other
without the underlying scale pattern itself changing first: a scaling
generator bracketed with \emph{itself} vanishes identically, and two
directions sharing a multiplet by definition share a scale value, hence
share a scaling generator, hence generate no internal rotation between
them — a binary fact about equal versus unequal scale values, not a
count that decreases gradually with multiplet size. Any change in which
directions belong to which multiplet accordingly requires a change in
the underlying scale pattern itself: multiplet structure cannot rotate
on its own, it can only be re-partitioned by the pattern that defines it
changing. What is unconditional, established directly in
Section~\ref{sec:particles} without this hypothesis, is that the
real-rank-one pattern's own rotation is carried entirely by its
variation, exactly as Section~\ref{sec:picture} states.

%======================================================================
\section{The isotropic sector: gravity}\label{sec:gravity}
%======================================================================

\subsection{The Schwarzschild solution and the two classical tests}

One honest cost of the isotropic sector is worth stating before its
results: over a commutative ring the general transport-commutator
construction of Section~\ref{sec:chain} gives \emph{no} curvature at
all, since the commutator that defines it vanishes identically on
commuting scalars — this is exactly why the isotropic geometry needs its
own, separately built curvature tensor rather than inheriting one from
the chain's general machinery, and it is a cost confined entirely to
this classical sector, not a limitation on the coupling itself, the
quantum sector, or the anisotropic sector, all of which keep their
transports genuinely non-commutative. With the scale pattern isotropic
in direction, a metric is built directly
from the scale field: writing the local scale along the time and radial
directions as $s_t, s_r$, the line element takes the diagonal form
$ds^2 = -s_t^2\,dt^2 + s_r^2\,dr^2 + r^2 d\Omega^2$. A deformation
tensor built directly from second derivatives of the scale field, in the
specific combination that the Riemann tensor is forced to equal for a
conformally flat metric, gives the Ricci tensor component by component
for this static, spherically symmetric ansatz; the
vacuum condition (Ricci flat away from a source) becomes the
stationarity of a particular sum of scale-derivative terms rather than
an independently imposed differential equation. Imposing asymptotic
flatness — $s_t, s_r \to 1$ at spatial infinity, the only boundary
condition consistent with A5 (no preferred finite reference scale) —
fixes the single remaining integration constant by itself, with no
further input: the stationarity condition only pins down the
\emph{derivative} of $\sigma_t+\sigma_r$, and asymptotic flatness alone
supplies the boundary value that makes the constant zero. A6$'$ plays no
role in this particular step; it enters later, in the Newtonian limit
and in fixing the overall normalization of the source term. Reciprocity
\begin{equation}
s_t \cdot s_r = 1
\end{equation}
follows as a theorem, not an ansatz: the Schwarzschild form is recovered
exactly, not approximately, as the unique solution consistent with the
axioms and the one external input A6$'$.

From $s_t s_r=1$, the PPN parameter is exactly $\gamma=1$, and the two
classical solar-system tests follow with no further assumption:
\begin{align}
\delta\phi_\odot &= 1.7515'' \quad(\text{measured } 1.7509\pm0.0002''\text{~\cite{Lambert2011}})\\
\Delta\varpi_{\text{Mercury}} &= 42.989''/\text{cent.}\quad(\text{measured } 42.98\pm0.04'')
\end{align}
Both bounds are certified to the precision quoted. Two internally
excluded alternatives are worth stating explicitly because they show the
framework can fail its own tests: computing deflection or precession
from the isotropic sector's magnitude change \emph{alone}, discarding the
directional structure of Section~\ref{sec:chain}, gives $0.8758''$ and
$14.33''$ respectively — both wrong by a clean integer factor (2 and 3)
relative to the measured values. Each is excluded by the same underlying
apparatus but through a distinct step: the deflection case is excluded
by the theorem that full isotropy together with reciprocity forces flat
space outright, while the precession case is excluded by the same
$\gamma=1$ derivation used for deflection, evaluated at the isotropic
sector's $\gamma=0$ instead. Neither exclusion is merely a comparison
with data.

\subsection{The Newtonian limit and gravitational redshift}

Linearizing $s_t = e^{\sigma}$ about the asymptotic value $\sigma=0$, the
scale response law A6$'$ reduces, to leading order in $\sigma$, to a
Poisson-type equation for $\sigma$ sourced by the same density $\rho$ that
appears in A6$'$, with the reciprocal constant fixed by the same
asymptotic-flatness argument used above. Since
$g_{tt}=-s_t^2=-e^{2\sigma}\approx -(1+2\sigma)$, a weak, static source
gives the standard Newtonian potential to leading order, recovering
Newtonian gravity as the linear regime of the same reciprocity relation
that gives $1.7515''$ at higher order — not a separately posited limit.
Gravitational redshift between two static observers at different $\sigma$
follows directly from the ratio of their local scale units,
$s_t(\sigma_2)/s_t(\sigma_1)$, and evaluated for a laboratory height
difference this reproduces the Pound--Rebka--Snider
result~\cite{PoundRebka1960,PoundSnider1965} to the precision of
Table~\ref{tab:scoreboard} — an arithmetic check of the A3 reading of
scale as a clock rate, not an independent prediction, since it uses the
same reciprocity relation already fixed above.

The linear regime above is not an approximation whose error one hopes is
small; it is realized \emph{exactly} by extending the substrate ring to
the dual numbers $A[\epsilon]/(\epsilon^2)$, a square-zero infinitesimal
extension, and setting $\sigma = \epsilon\varphi$ for an ordinary
function $\varphi$. Every derivation carries over componentwise to this
extended ring, and the quadratic invariant $|\nabla\sigma|^2$ vanishes
\emph{identically} rather than merely being small, since
$\sigma_i^2 = \epsilon^2(\partial_i\varphi)^2 = 0$ by $\epsilon^2=0$.
Poisson's equation $2(n-1)\Delta\Phi = \kappa\rho$ therefore holds with
no term discarded and no linearization error, for the entirely
algebraic reason that there is nothing to discard: the ``weak-field
limit'' is not a truncation of the exact theory but an exact statement
about a different, infinitesimally-deformed ring, related to the finite
theory only by which numbers one is willing to admit as the value of
$\sigma$.

\subsection{Local covariance survives global dissipation}

A6 makes the theory globally non-conservative: there is no functional
whose gradient reproduces the total drift of the scale field. This does
not threaten local Lorentz covariance, because the two facts live at
different scopes. The coupling of Theorem~\ref{thm:coupling} and the
Lorentz algebra of Section~\ref{sec:chain} are pointwise, algebraic
facts, holding identically at every point regardless of how the scale
field as a whole evolves; A6 is a global, integrated statement about the
\emph{sum} of the drift across all points. A local Lorentz frame is
constructed at each point from the coupling alone and is exactly as
covariant there as in a theory with a globally conserved potential;
openness only says that no single potential accounts for the drift
\emph{everywhere at once}, which is compatible with, and does not
weaken, covariance at any one point.

\subsection{No singularity, no horizon}

Because $g_{tt}=-s_t^2$ with $s_t$ a unit — nonzero everywhere by A2 —
$g_{tt}$ never vanishes: there is no horizon and no curvature
singularity. The exterior geometry away from the would-be horizon is
unchanged from Schwarzschild, so shadow shape, ISCO location, and the
early ringdown spectrum are unaffected; the difference between this
framework and general relativity is concentrated entirely in the
boundary condition at small radius, where general relativity imposes
absorption (a horizon) and this framework cannot, since nothing in it is
ever exactly zero.

The absence of a singularity is not merely the absence of a place where
$s_t$ hits zero; it is structural at the level of what curvature is
built from. The Riemann tensor of Section~\ref{sec:chain} is a fixed
linear function of a deformation tensor built entirely from the
\emph{first and second derivatives} of $\sigma$ — the scale value $s$
itself never appears in the curvature formula, only its rate of change.
Two scale fields with identical derivatives at every point therefore
have identical curvature even if their values differ arbitrarily, and a
region where the scale becomes extreme is on that account alone no more
curved than a region where it does not: what general relativity would
read as an approach to a singularity, this framework reads as a region
where $\sigma$ is large in magnitude but changing no faster than
elsewhere, which is exactly what asymptotic flatness combined with
reciprocity $s_ts_r=1$ guarantees, since a large $\sigma$ in one
direction is bought by an equally large and opposite $\sigma$ in the
reciprocal direction, without either derivative diverging.

A scale ratio, being built from units, is itself a unit — never zero,
never infinite, but with no upper bound on its magnitude either. There
is consequently no minimum resolvable length and, on the matter side of
the same statement, no maximum density
(Section~\ref{ss:native}): for any proposed shortest scale a finer one
is always available, so there is no final effective description this
framework's own structure ever forces a halt at, only successive
replacement at each threshold crossed.

\subsection{Two polarizations, from two different sectors, and no breathing mode}\label{ss:pol}

A transverse gravitational wave in $k$ spatial dimensions can, a priori,
carry perturbations from two independent sources at each point on the
propagation plane: a change in the local scale magnitude (the scaling
sector) and a change in the local scale direction (the rotation sector),
matching the split of Section~\ref{sec:chain}. In the two transverse
directions available for $k=3$, the scaling sector contributes
$\dim(\text{transverse}) - 1 = 1$ independent mode — not $2$, because a
uniform rescaling of both transverse directions together is exactly a
fiducial shift, and A5 declares fiducial shifts unobservable, removing
that mode from the physical spectrum entirely rather than merely making
it small. The rotation sector contributes $\dim\Lambda^2(\text{transverse}) = 1$
independent mode, the wedge of the two transverse directions. The total
is
\begin{equation}
1_{\text{scale}} + 1_{\text{rotation}} = 2,
\end{equation}
matching the two tensor polarizations of gravitational-wave observation,
but arrived at as a sum of contributions from two structurally distinct
sectors rather than as a single count. This gives a discriminator sharper
than the count alone: the framework has \textbf{no scalar (breathing)
mode at any amplitude}, since that mode is precisely the uniform
rescaling A5 removes, whereas scalar–tensor theories of gravity generically
predict one at some coupling-dependent amplitude. A breathing-mode
detection would not merely disfavor this framework; A5 gives it nowhere
to appear at all, so the correct statement is not ``predicted small'' but
``predicted absent.'' This also bears, incidentally, on the ultralight-boson
exclusion drawn from GW241011's rapid spin~\cite{GW241011}: that analysis
rules out a scalar of Compton wavelength in the superradiant band, and this
framework predicts no such field to begin with, since its label structure
(threshold, $\mathbb{Z}/2$, $\mathbb{Z}$) has no further slot for one and the
scalar polarization it would carry is exactly what A5 removes — the
non-detection needs no explanation here, where a scalar-tensor theory owes one.

\subsection{Gravitational waves: what radiates, and what does not}

Conservation is not an additional postulate laid alongside a field
equation for the scale field; it is the condition for such an equation
to be writable at all. If curvature is proportional to a source — the
shape of any field equation — the Bianchi identity forces that source's
cyclic covariant divergence to vanish identically, with no further
input, since Bianchi is itself a consequence of the associativity of
transport (the Jacobi identity for the transport bracket). The
contrapositive is the sharp form: a source whose cyclic divergence fails
to vanish cannot be the curvature source of \emph{any} connection, so
this framework does not merely prefer conserved sources over
non-conserved ones — it has no expression for a non-conserved source at
all. Conservation, read this way from the Bianchi identity as a
statement about the scale field's own algebra, kills the monopole and
dipole moments of radiation specifically. This is established purely algebraically: divergences form the
additive subgroup generated by images of the non-drift derivations, and
because A1's derivations commute, the drift direction descends
consistently to the quotient by that subgroup — no measure, no Stokes'
theorem, and no decay condition at infinity is used; what survives is one
purely algebraic condition, that this subgroup of divergences is proper
rather than the whole ring. Dipole radiation is accordingly identically
zero, consistent with the double-pulsar bound
$|\dot P_b^{\text{excess}}/\dot P_b|<1.3\times10^{-4}$; the quadrupole
leads, giving the standard inspiral chirp.

The one constant $\kappa$ supplied by A6$'$ appears exactly once, not
twice: it fixes both the Newtonian-limit Poisson normalization and the
quadrupole radiation coefficient, so fixing it from the static, weak-field
limit leaves the radiated power with no remaining freedom — a binary
pulsar's orbital decay is consequently a genuine test rather than a fit
to a second free constant. For a circular binary with component masses
$m_1, m_2$ and separation $r$, the orbital energy $E=-Gm_1m_2/2r$ and the
quadrupole radiated power
$P = \tfrac{32}{5}G^4(m_1m_2)^2(m_1+m_2)/(c^5r^5)$ together force the
orbit to shrink, $dr/dt<0$, at a closed-form rate obtained by nothing
more than the chain rule applied to $dE/dt=-P$. The gravitational-wave
frequency evolution depends on the two masses only through the chirp
mass $\mathcal{M}=(m_1m_2)^{3/5}/(m_1+m_2)^{1/5}$, via
$df/dt \propto \mathcal{M}^{5/3}f^{11/3}$ — the same combination and the
same two exponents, $5/3$ and $11/3$, that a gravitational-wave detector
template actually fits. The chirp mass combination and its two exponents
are read off directly from the same quadrupole formula used for the
binary-pulsar prediction; the further step from the time-domain
separation rate to this frequency-domain form uses Kepler's third law in
its standard, uncontested form, which is cited rather than re-derived
here — the exponents are not an independent input chosen to match
observation, but the derivation is not carried through in full from
first principles either. Applied to PSR~B1913+16
(Hulse–Taylor)~\cite{HulseTaylor1975,WeisbergHuang2016}, the
predicted-to-observed
ratio of orbital decay is $0.9983$, agreeing to two parts in a thousand
and bounding any dipole contamination below $0.2\%$; applied to the
double pulsar PSR~J0737$-$3039~\cite{Kramer2021}, sixteen years of timing bound the same
dipole fraction below $1.3\times10^{-4}$, an order of magnitude tighter,
with every improvement in timing precision tightening a prediction that
has no parameter left to retune.

\begin{theorem}[All-order equivalence with general relativity]\label{thm:allorder}
The symmetric (radiative) part of the deformation tensor equals its
classical, general-relativistic counterpart at \emph{every}
post-Newtonian order.
\end{theorem}

This is an exact statement, not a leading-order approximation, and its
origin is given in Section~\ref{sec:quantum}: the entire quantum
correction to curvature is purely antisymmetric, so the symmetric sector
carries none of it, at any order. There is accordingly no order at which
the two theories' gravitational-wave predictions could differ — not
merely an order too small to currently measure. The gravitational-wave
sector can discriminate this framework from general relativity only
through the ringdown boundary condition, the presence or absence of a
horizon, never through the inspiral or the post-Newtonian expansion
itself.

\subsection{The horizon-absence signature: echoes and the ergoregion bound}\label{ss:echo}

Absence of a horizon permits, but does not predict the amplitude of, late
echoes. Writing $A_{\min}$ for the field's reflectivity at the would-be
horizon (a quantity this framework does not compute, Section~\ref{ss:eos-out})
and $\tau$ for the light-travel time to that surface, the round-trip
delay is $\tau_{\text{echo}}=\tau\ln(1/A_{\min})$ — logarithmic, not
linear, in the redshift $1/A_{\min}$: ten orders of magnitude of
additional redshift lengthen the delay by only a factor of order ten, not
$10^{10}$, which is why an echo search looks for a comb of arrivals at
millisecond spacing rather than at a Planck time. The $n$-th echo arrives
at $n\,\tau_{\text{echo}}$, evenly spaced, with amplitude falling as
$A_0 R^{2n}$ for reflectivity $R$ after $n$ round trips — a null search
excluding an amplitude above some fraction $b$ of the ringdown directly
bounds $R^2 \leq b$.

Two independent observational handles exist, and they scale with $R$
differently. Echo searches place a direct bound on $R$ that is
polynomial in the excluded amplitude fraction, as above. Ergoregion
(superradiant) instability of a rapidly spinning, long-lived compact
remnant places an \emph{exponential} bound on $\mathcal{R}$, because a
horizonless spinning object whose reflectivity is not exponentially
close to unity is unstable on an observationally short timescale — the
growth rate of the instability, not merely the detector's sensitivity to
a faint echo, sets the bound, so a longer-lived remnant tightens the
constraint exponentially rather than linearly. GW241011's~\cite{GW241011}
primary spin, $\chi_1=0.64>0.55$ at the reported $90\%$ lower limit, is in the rapid
regime, and the remnant's evident long-term stability makes this the
tighter of the two constraints currently available — and, honestly
stated, a real and previously unregistered pressure on this framework:
because $\mathcal{R}$ is not computable here (Section~\ref{ss:eos-out}),
the framework cannot currently say whether it survives this constraint,
only that the constraint applies and tightens with every additional
rapidly-spinning long-lived remnant observed. The comparison
between GW250114's~\cite{GW250114} ringdown and its inspiral-predicted remnant is itself
a genuine confirmation, independent of the reflectivity question: the
inferred final mass ($62.7\,M_\odot$, from components $33.6$ and
$32.2\,M_\odot$ radiating $3.1\,M_\odot$, about $5\%$ of the total) and
final spin ($0.68$) predict a ringdown spectrum whose first overtone
matches the fundamental-mode inference to within the Kerr-family bound
of $30\%$, and the area law — that the horizon area cannot decrease — is
confirmed at $3.4\sigma$ across every analysis window and above $5\sigma$
in the most sensitive one; both are area-law-family statements this
framework's exterior geometry inherits unchanged from general relativity,
since nothing about the exterior solution differs. The damping-time
precision of that same ringdown measurement ($9\%$) is looser than the
frequency precision ($2\%$) by a factor of about $4.5$, which matters because
the damping time is exactly the observable a partially reflecting
boundary would shift, while the frequency would not — so the ringdown
measurement is currently most sensitive in precisely the channel where a
reflecting-boundary effect would first appear, and has not yet reached
the precision needed to see one.

The magnitude of $\mathcal{R}$ itself requires an
equation of state for whatever replaces the interior — a computation
explicitly outside this framework's scope (Section~\ref{ss:eos-out}) —
but the qualitative structure, comb-like echoes and an exponential rather
than linear bound, is a genuine, falsifiable prediction independent of
that missing input.

%======================================================================
\section{The anisotropic sector: labels and charges}\label{sec:particles}
%======================================================================

Where the scale pattern is anisotropic, an order parameter, charges, and
multiplets appear together, from the same source as the gauge connection
of Section~\ref{sec:chain}. A particle is a point defect in the scale
vector field — an unresolvable knot — completely characterized by a
triple: a threshold (its position on the scale axis, read as mass), a
$\mathbb{Z}/2$ confined charge, and a $\mathbb{Z}$ winding charge, and
nothing else. This is not an assumption about how many labels a particle
may carry; the pattern having collapsed to a single vector
(Section~\ref{sec:chain}) leaves two topological invariants available in
three spatial dimensions — a winding number (valued in $\mathbb{Z}$,
from the vector's degree as a map to a sphere) and a monodromy class of
the order parameter (valued in $\mathbb{Z}/2$, because the order
parameter lives in a projective, not a linear, space). That these are
the invariants a projective order parameter in three dimensions admits
is a standard fact of algebraic topology, cited rather than proved
in-repo; what \emph{is} proved directly is the sharper, load-bearing
half — that the confinement group is specifically $\mathbb{Z}/2$ and not
some other value, which is what rules out the three-valued color
identification below.

The $\mathbb{Z}/2$ nature of the confined charge is structural rather
than a numerical near-miss: an earlier attempt to read this charge as
color required three distinguishable values, and failed against the data
by giving two. The order parameter's home is a projective space, and a
projective order parameter simply does not have a three-valued invariant
available to it — no adjustment of the construction changes this.

Mass, in this picture, is not a third label alongside the two charges; it
is the position on the scale axis at which unresolved content first
becomes a distinguishable defect, so ``mass'' and ``threshold'' are the
same statement read two ways rather than an independent property a
particle additionally carries. This is also the precise sense in which
there is no final, complete list of fundamental particles
(Section~\ref{sec:picture}): resolving further along the scale axis
always exposes more threshold structure above whatever has already been
resolved, and the claim is not merely that undiscovered particles might
exist but that, given an unbounded supply of ever-finer thresholds to
resolve — itself a physical premise about the scale axis rather than a
theorem derived from A1–A7 alone — the content strictly above any given
resolution is never empty, so ``the complete particle table'' is not a
question this framework's content structure can answer even in
principle, for any resolution one names.

\subsection{The two charges are independent, and exactly one is confined}

The winding charge and the monodromy (block) charge are logically
independent: neither is determined by the other, since they are computed
from different topological data — a degree and a monodromy class — of
the same underlying configuration, and it is possible to exhibit
configurations agreeing on one label while differing on the other. This
independence is what makes ``particle'' a genuine pair of quantum
numbers rather than one label in disguise, and it is what makes the
statement ``exactly one of the two is confined'' non-trivial rather than
automatic: confinement is a further, structural fact about the
monodromy charge specifically — a configuration whose block labelling
fails to close up after a circuit is not merely rare or energetically
disfavored, it is not single-valued and so does not describe a state
that can exist in isolation — while the winding charge carries no such
obstruction and every winding number is freely realizable on its own. A
defect paired with its own opposite always closes up and carries zero
net charge in \emph{both} labels simultaneously, independent of which
block permutation either one carries, which is the framework's account
of why bound pairs (mesons, in the language of the withdrawn color
identification) are always available regardless of the confined
partner's internal labelling.

The same rank-one collapse that limits particles to two labels also
bounds the shape of the mass spectrum. If particle masses were drawn from
a mixture of two or more structurally distinct populations rather than a
single rank-one family, the squared coefficient of variation of the mass
distribution satisfies $\mathrm{CV}^2 \geq 1$, with equality exactly when
the population is the single-family (rank-one) case this framework
predicts; a mixture is not merely disfavored numerically, it is excluded
whenever the measured $\mathrm{CV}^2$ is below $1$. The measured value,
$0.875$ (Table~\ref{tab:scoreboard}), sits below the bound and is
consistent with rank one, though $0.875 \neq 1$ exactly, so the
comparison establishes exclusion of mixtures rather than confirmation of
the exact rank-one value — a distinction Section~\ref{sec:register}
records explicitly rather than eliding.

\subsection{Gauge symmetry is the stabilizer of the pattern}

Internal symmetry in this framework has no independent space to act on;
the only structure at a point is the scale pattern across directions, so
any symmetry that is not a motion of the substrate itself must be a
relabeling of directions that leaves every local scale value unchanged.
The group of such relabelings — the \emph{stabilizer} of the pattern
under the permutation action on directions — is not chosen; it is read
off the pattern directly, and two directions belong to the same
multiplet exactly when they carry the same local scale, so multiplets
\emph{are} the degeneracy classes of the pattern and multiplet sizes are
the sizes of those classes. Where the pattern is injective (no two
directions share a value) the stabilizer is trivial and there is no
internal symmetry at all; where the pattern is fully isotropic the
stabilizer is the entire permutation group.

This has a direct consequence for what a connection is. A transport
between two scale patterns is a relabeling that matches them; where the
target pattern is injective, that relabeling is unique — the connection
is a function of the scale field alone, with nothing left to choose —
and wherever it is not unique, the freedom in choosing it is exactly a
coset of the stabilizer, meaning gauge freedom is not a separate
structure layered onto the geometry but is \emph{identical to} the
degeneracy of the underlying pattern: gauge freedom exists at a point if
and only if that point's scale pattern is degenerate there, and its size
is exactly the size of the degeneracy. A theory with no coincident scale
values anywhere has no gauge freedom anywhere, not as a choice of gauge
but as a consequence of the pattern being injective; internal symmetry is
consequently never posited as an independent input; it is measured by
counting coincidences in a single scalar-valued function per direction.

\subsection{Particle lifetimes are scale ratios}

Because energy scale and length scale are reciprocal (A3), a particle's
proper lifetime $\tau_0$ and its lab-frame lifetime $\tau$ are related by
the same ratio that appears everywhere else in this theory:
\begin{equation}
\tau = \gamma\,\tau_0, \qquad \gamma = \frac{\varepsilon}{\varepsilon_0},
\end{equation}
the ratio of the particle's energy scale in the lab frame to its energy
scale at rest. Time dilation is accordingly not a separate relativistic
postulate but the same scale-ratio statement that produces every other
result in this paper, evaluated for a moving particle; applied to the
CERN muon storage ring measurement~\cite{Bailey1977} this reproduces
the measured lifetime to the precision of
Table~\ref{tab:scoreboard}.

\subsection{Is there a second switch?}\label{ss:secondswitch}

The framework predicts exactly one kind of label-opening event: an
anisotropic pattern appearing where an isotropic one was, always in the
direction of \emph{refinement} — a multiplet only ever splits into
smaller pieces as the scale is probed, never the reverse. The framework
does not itself prove this refinement direction — only that a change of
label structure requires a change of scale pattern — so testing the
refinement direction against data is not circular.

We count exactly-conserved label \emph{types} (not states) across the
directly probed particle spectrum~\cite{PDG2024}, spanning
\begin{equation}
\frac{E_{\max}}{E_{\min}} = \frac{1.3\times10^{4}\,\text{GeV}}{5\times10^{-11}\,\text{GeV}} > 10^{14},
\end{equation}
fourteen orders of magnitude in energy. Over that range the electroweak
multiplet only ever shrinks ($2\to1$, a strict refinement) and the color
multiplet is unchanged ($3\to3$); no observed symmetry breaking has ever
enlarged a multiplet, and no additional exactly-conserved label beyond
the four already in the framework's inventory has been found. This is a
genuine, if limited, empirical check: it can only detect \emph{exact}
degeneracies opening or closing, and fourteen decades, while large, is
finite.

%======================================================================
\section{Scale flow: renormalization group and the strong scale}\label{sec:flow}
%======================================================================

A fiducial shift is a symmetry of the geometry (A5) but not automatically
of the description of matter: to keep predictions fixed while the
reference unit moves, the couplings used to describe matter must move to
compensate, and that compensating motion is the renormalization group.
Three structural facts follow from A5 alone, before any coupling is
specified. First, the scale flow is forced to be a one-parameter group —
shifting the fiducial by $t$ and then by $u$ is shifting it once by
$t+u$, with no other composition law available. Second, self-similarity
forces exponential solutions: any observable obeying the
Callan–Symanzik-type equation $dO/dt=\gamma O$ satisfies
$O(t)=O(0)e^{\gamma t}$ exactly, the same mathematical fact that produces
exact de Sitter expansion in Section~\ref{sec:cosmology} — the same
equation read at two different scales, not the same numerical constant
in both places. Third, the effective-field-theory tower is forced: each
excitation carries a log-threshold, the set of excitations available at
a given log-energy is monotone increasing in that energy, and it is
exactly constant between thresholds — the description does not merely
happen to be piecewise constant, it is provably unchanged wherever no
threshold is crossed. These three facts fix the \emph{shape} of any beta
function forced by A5: piecewise constant, jumping only where the
active content changes. They do not, by themselves, fix the coefficient
at any single jump — that remains external input at this level of the
argument.

The strong-interaction sector closes that remaining gap for one specific
coupling. Reading a coupling not as a diagrammatic sum but as an
accumulated \emph{count} — how many resolvable defects a probe at a
given log-scale has responded to — A5 forces that count's density to be
constant in log-scale, since a scale-covariant world has no preferred
log-scale for defects to bunch around. A constant density $\rho$
integrated over an interval is linear, giving
\begin{equation}
g^{-2}(t) = g^{-2}(t_0) + \kappa\rho\,(t-t_0),
\end{equation}
exactly the one-loop running form, with the one-loop coefficient
identified as $b=\kappa\rho/2$ — a density of defects per unit log-scale
in place of an imported group-theory coefficient (with $\kappa>0$
retained as a hypothesis, not itself derived). Given that, asymptotic
freedom is \emph{equivalent} to $\rho>0$: $g(t)^2\to0$ as $t\to\infty$
exactly when there are defects to resolve, so which sign of density
gives a free theory is read off the count rather than fitted
separately. The combination
$\Lambda := \exp\bigl(t-g(t)^{-2}/(2b)\bigr)$ is constant along this
flow — a scale-\emph{invariant} axiom system has produced a definite
scale not by breaking the invariance in the axioms but by the solution
of a scale-covariant equation secreting one, and A5 guarantees no choice
of fiducial can shift it. Every hadron mass, read as $c\cdot\Lambda$ for
a pure number $c$, then has a mass \emph{ratio} $c_1/c_2$ independent of
$\Lambda$ and hence of every free input to the theory, in contrast to
the overall scale $\Lambda$ itself, which is measured rather than
predicted, consistent with the one-free-unit statement running through
this paper.

One caution is necessary about $\rho$, and it is the framework's own
internal correction rather than an external criticism. The density
entering the running above is \emph{sector-restricted and signed}: a
coupling responds only to the structures actually charged under it, not
to every threshold in the spectrum, so it is emphatically not the same
quantity as the raw threshold density used elsewhere
(Section~\ref{sec:register}) to bound the mass spectrum. A proposed
cross-check — measuring this running-coupling density independently from
the raw spectral count and comparing — was attempted and found not to be
a test at all: the two counts differ by roughly a factor of two ($0.864$
counting all twelve Standard Model mass scales against $0.444$ counting
the six quarks alone), and within a single sector the relating equation
has one free coefficient and one measurement, so it fixes the
coefficient rather than checking anything. This is registered
(Section~\ref{sec:register}) as a proposal that was made and retracted,
not as a claim that survives.

Measured against data, the simplest possible ratio law for the confined
sector — an equal-ratio geometric tower of defect masses — was checked
against lattice glueball ratios~\cite{MorningstarPeardon1999} ($1.39$–$1.40$ and $1.07$–$1.08$ for the
two lightest ratios, differing by more than $0.3$) and
\textbf{failed}, refuting that specific hypothesis about the mass law
rather than the parameter-free ratio construction itself
(Section~\ref{sec:register} records the retraction and why the
underlying flow result survives it).

\subsection{A second, sharper test: the shape of the threshold spectrum itself}

The glueball comparison tests one specific alternative hypothesis using
three numbers. A5's requirement that no log-scale be preferred makes a
second, independent, and considerably sharper prediction using the
entire spectrum at once. If thresholds are placed by a scale-invariant
process — one of constant rate $\rho$ per unit log-scale, with no
preferred location for that rate to be higher or lower — the resulting
point process has exponentially distributed gaps, and an exponential
distribution has squared coefficient of variation exactly
$\mathrm{CV}^2=1$, with no free parameter. The geometric-tower hypothesis
predicts the opposite extreme: equal gaps, hence zero variance, hence
$\mathrm{CV}^2=0$ exactly. The two predictions are as far apart as they
can be, and this is a statement about the whole gap distribution, not
about two or three selected ratios.

Taking the eleven consecutive gaps in $\ln m$ across the twelve distinct
Standard Model mass scales, from the electron to the top quark, gives
$\mathrm{CV}^2=0.875$ ($\mathrm{CV}=0.935$) — inside the central $90\%$
sampling range for eleven gaps under the scale-invariant prediction
(roughly $0.33$ to $1.64$ in $\mathrm{CV}^2$), and at the $63$rd
percentile of it, while the geometric tower's $\mathrm{CV}^2=0$ lies far
outside that range on the other side entirely. What this establishes,
honestly stated, is the sharper half of the comparison: exclusion of the
geometric tower is decisive, since $0.875$ is nowhere near $0$, while
agreement with the scale-invariant value $1$ is merely encouraging,
since eleven gaps is a small sample and the twelve mass scales entering
it were selected by hand (massless states excluded, neutrino masses
unresolved). The same gap statistics fix the threshold density,
$\rho\approx0.86$ per e-fold — the same $\rho_{\text{all}}$ that appears
in the retracted cross-sector comparison above, now correctly understood
as a property of the full spectrum's gap statistics rather than a
quantity meant to match a sector-restricted running coefficient.

\subsection{Label proliferation measures range, not structure}\label{ss:labelrange}

A description of content is exact between thresholds, as noted above,
and gains exactly the thresholds it
crosses when its range is widened, and no others: extending the window
of validity from $t$ to $t'$ adds precisely the structures whose
threshold lies in $(t,t']$, and if no threshold lies in that interval
the description is unchanged. Two consequences follow directly, with no
further assumption. First, the excess content a formalism must carry is
monotone in the range it covers — widening the window at either end can
only add labels, never remove them, so the label count is a genuine
function of how much log-range is covered. Second, no finite label set
suffices at every scale: above any given scale there is always a
threshold not yet crossed, so any fixed labeling scheme is eventually
insufficient, not as a practical limitation but as a theorem.

This gives a direct account of why the Standard Model's label set is as
large as it is, without touching whether its content is fundamental. A
formalism required to remain valid across thirty orders of magnitude in
energy must carry a label for everything that appears \emph{anywhere} in
that range, whether or not it is realized at the particular scale one
happens to probe — most of those labels are, at any single scale,
unrealized. This is exactly the shape of off-shell content, gauge
redundancy, and species that exist only virtually: not evidence that the
underlying structure is as rich as the label count suggests, but the
unavoidable bookkeeping cost of a single formalism spanning a wide
range. The claim is not that the Standard Model is wrong to carry the
labels it carries, only that the size of that label set is telling an
observer about the range being covered, and reproducing that particular
count is consequently not a target this framework treats as meaningful.

%======================================================================
\section{The quantum sector}\label{sec:quantum}
%======================================================================

Promoting the commutative substrate to a non-commutative one and asking
what changes gives the quantum sector. The commutators of the
non-commutative scale algebra reproduce the Leibniz rule, the Jacobi
identity, the power rule for derivatives, and the canonical commutation
relation as theorems rather than postulates. Consequently, A3's scale-energy
duality $\varepsilon\cdot s=1$ continues to hold exactly only up to a
correction of order $\hbar$ in the non-commutative sector — the classical
relation is the $\hbar\to0$ limit of a relation that is otherwise
corrected, not an independent low-energy postulate.

Uncertainty is a general fact about symmetric (self-adjoint) operators on
an inner-product space and does not require the scale algebra
specifically: for any two symmetric operators $A,B$ and any unit vector
$\psi$,
\begin{equation}
|\langle\psi,[A,B]\psi\rangle| \le 2\|A\psi\|\,\|B\psi\|,
\end{equation}
by the Hermiticity of $A,B$, the triangle inequality, and Cauchy–Schwarz
applied to the resulting inner products. Specializing to a commutator
that acts as the scalar $i\hbar$ — which is what the canonical
Poisson-bracket construction below produces for position and its
conjugate momentum, once quantized — gives the Heisenberg relation
$\Delta A\cdot\Delta B \ge \hbar/2$ as a corollary of the general
inequality, not as an independently assumed postulate: the constant
$\hbar$ enters only through evaluating the commutator, and the
inequality itself is pure linear algebra.

Before the exact (Theorem~\ref{thm:hbar}) construction, a first-order
deformation-quantization~\cite{Bayen1978} version fixes what
quantization must mean here. On $A[\hbar]/(\hbar^2)$,
a square-zero extension of the substrate ring, a star product is built
directly from a biderivation $P$ (a Leibniz bilinear pairing of two
derivations, not itself required to be antisymmetric):
$(f_0,f_1)\star(g_0,g_1) := (f_0g_0,\, f_0g_1+f_1g_0+P(f_0,g_0))$. This
product is associative and reduces exactly — with nothing approximated —
to the ordinary commutative product at $\hbar^0$, and its commutator
reproduces the \emph{antisymmetric part} of $P$ at $\hbar^1$:
$\hbar^{-1}[f,g]_\star = \{f_0,g_0\}_P := P(f_0,g_0)-P(g_0,f_0)$, the
Poisson bracket, so two quantities commute in the deformed algebra
exactly when this antisymmetric part vanishes on them, and the
commutative (gravitational) sector is recovered exactly as the locus
where that bracket is zero — $P$ itself need not vanish there, only its
antisymmetrization. This is a first-order formal deformation, useful for
identifying what non-commutativity must mean
before asking whether it can be realized exactly; Theorem~\ref{thm:hbar}
below removes the formal parameter and replaces it with an honest
additive-and-multiplicative shift map.

\begin{theorem}[$\hbar$ as an exact scale step]\label{thm:hbar}
There exists a scale shift $T$ — a map on the scale algebra additive and
multiplicative, though not required to fix the multiplicative identity —
such that $T$ fails to commute with position by exactly one unit:
$\Delta X := T(X)-X=1$, exactly, not as a leading term in an expansion.
\end{theorem}

This structure is never itself constructed from A1–A7 alone by the
argument above; it asks only for a map with this weaker-than-ring-endomorphism
property.
We therefore exhibit an explicit model: the polynomial ring
$\mathbb{R}[X]$ with $d/dX$ satisfies A1 for $n=1$, and the translation
$X\mapsto X+1$ is a ring homomorphism realizing $T$ exactly, so the shift
is a legitimate, exact, finite — not infinitesimal — translation along
the scale axis. This shows the theorem is non-vacuous and the structure
compatible with the axioms; it does \textbf{not} show that every model of
the axioms must carry such a shift, so the identification of $\hbar$ with
this step remains an identification, not a derivation, and is registered
as such (Section~\ref{sec:register}).

\begin{theorem}[The quantum correction to curvature is purely antisymmetric]\label{thm:ncconf}
The entire non-commutative correction to the deformation tensor is given
by the commutator $[\sigma_i,\sigma_j]$ of scale generators, which is
purely antisymmetric.
\end{theorem}

This is the origin of Theorem~\ref{thm:allorder}: since the symmetric
sector carries no quantum correction at all, the gravitational-wave
radiative sector is untouched by non-commutativity to any order, while
the antisymmetric (rotational/gauge) sector carries the entire quantum
correction. This gives a genuine, structural prediction for \emph{where}
a quantum-gravitational effect could appear at all, even though its
magnitude cannot be: the correction $[\sigma_i,\sigma_j]$ and the
rotational label of Section~\ref{sec:chain} are the same kind of
object, both antisymmetric in the same index pair, so
\textbf{quantum-gravitational corrections appear first in
spin–geometry coupling, not in orbital dynamics} — orbits are computed
from the symmetric part, which is exactly the part this correction
cannot touch. Stated honestly rather than left implicit: the size of
the commutator is set by the same scale step that gives $\hbar$
(Section~\ref{sec:quantum}), so relative to the classical terms in the
solar system this correction is utterly negligible, and nothing here
predicts an effect within reach of current instruments — what is
claimed is the structure of the leading correction and precisely which
channel it can and cannot appear in, not a measurable number.

%======================================================================
\section{Dark matter, dark energy, and light}\label{sec:cosmology}
%======================================================================

\subsection{Dark matter: two routes, making opposite predictions}

This framework offers not one account of dark matter but two, and it
does not choose between them — an honest state of affairs rather than an
oversight, since the two are logically independent and observationally
distinguishable.

The first is a detection-threshold mechanism, requiring no branch-point
argument at all. A detector is itself a physical system sitting at some
energy scale $\varepsilon_D$; by A3 that scale is the reciprocal of its
spatial resolution, so what it can register is not ``an interaction''
but an interaction that moves its own state by at least its threshold
$\theta$. An excitation whose own energy scale $\varepsilon$ sits far
below $\varepsilon_D$ deposits a fraction $\varepsilon/\varepsilon_D$ of
the detector's scale — strictly positive, so the interaction genuinely
happens, but driven below threshold whenever $\varepsilon$ is small
enough, so the observation genuinely fails. Gravity, meanwhile, couples
to the total energy $N\varepsilon$ of a population of such excitations,
which is blind to how that total is partitioned among them. A large
multiplicity of very-low-scale excitations is consequently
gravitationally loud and non-gravitationally silent \emph{simultaneously
and for the same population}, for any fixed gravitational signal and any
fixed detector: the theorem constructs, for any target mass $M$ and any
detector threshold, a scale and multiplicity reproducing $M$
gravitationally while sitting strictly below that threshold. On this
route dark matter has thresholds and hence a genuine mass spectrum, and
a sufficiently well-matched probe would eventually see a signal.

The second is the isotropic-trajectory mechanism of
Section~\ref{sec:branch}: a region whose scale pattern is isotropic in
direction still has curvature, since magnitude inhomogeneity, not
direction inhomogeneity, sources gravity, so such a region gravitates
with no anisotropic order parameter and hence no electromagnetic or
confined charge (Section~\ref{sec:particles}) at any scale whatsoever —
not a small coupling waiting to be resolved, but no label to couple to
in the first place. On this route dark matter has no thresholds and
hence no mass spectrum and no species to enumerate: it is one structural
condition on the scale pattern, not a family of candidate particles with
different masses and couplings, and it forbids exactly what the first
route permits — a charge-type self-interaction, since that requires the
directional order parameter this route's matter provably lacks.

The two routes make \textbf{opposite} predictions and the framework
states this as a fork rather than resolving it: the low-scale route
gives a coupling that is small but strictly nonzero at every finite
separation, hence a spectrum in principle detectable by a sufficiently
sensitive or well-matched probe; the isotropic route gives a rotational
label that is identically zero at \emph{every} scale, hence collisionless
behavior and no spectrum to find by any amount of improved sensitivity.
A null result at ever-improving direct-detection sensitivity favors the
isotropic route; a spectral feature or a self-interaction signature
favors the low-scale route; the framework hosts both constructions and
does not adjudicate between them, which is a weaker but more honest
position than committing to either.

\subsection{The direction of time from finite resolution alone}

The entropy ordering used in Theorem~\ref{thm:onearrow} rests on an
$H$-theorem that uses nothing about scale in its own proof, which is
precisely what makes the later identification with the scale-drift
arrow a substantive fact rather than a definition chasing its own tail.
The construction is elementary: let $X$ be a finite set of
microscopic configurations and $\pi:X\to Y$ a resolution map, generally
non-injective because observation happens at finite resolution. The
\emph{saturation} of a subset $S\subseteq X$ under $\pi$ is the set of
all configurations $\pi$ cannot distinguish from something in $S$; a
configuration set is always contained in its own saturation, so
coarse-graining can only add configurations, never remove them, giving
$|S|\le|\mathrm{sat}(\pi,S)|$. Microscopic evolution, being a bijection
(the substrate is reversible — nothing in A1–A7 privileges a direction of
time at the microscopic level), preserves the cardinality of any subset
exactly. Composing one step of reversible evolution with one step of
coarse-graining therefore cannot decrease the number of configurations
compatible with what is observed: entropy, read as the log of that
count, is non-decreasing by counting alone, with no statistical
assumption beyond finite resolution — and once it has strictly increased
at one step, no entropy-non-decreasing law can undo the evolution,
which is the precise sense in which a reversible microscopic law
produces an irreversible macroscopic one, matching the status of the
Boltzmann equation.

This also gives the framework's structural reading of Sakharov's third
condition for generating any asymmetry at all: departure from thermal
equilibrium. In a closed system that departure is a fine-tuned episode
that must be arranged; under A6 the universe never stops dissipating,
so the condition is satisfied permanently rather than accidentally —
strict entropy production is itself a departure from equilibrium, since
every observable is stationary in equilibrium by definition, so an open,
dissipating universe satisfies Sakharov's condition at every moment it
produces entropy at all. This is deliberately modest: it establishes
that the necessary condition holds throughout this framework's history,
not that any particular observed asymmetry (baryon or otherwise) follows
from it.

\subsection{Dark energy is the dissipation rate, squared}

Cosmic expansion is, in this framework, not the motion of anything
through space but the running down of the total dimensionless energy
scale $\varepsilon$, read through the duality A3 as a growing length
scale. Dissipation at a constant fractional rate $\lambda$ has a unique
solution, $\varepsilon(t)=\varepsilon_0e^{-\lambda t}$ — no ansatz is
assumed, the exponential is forced by the differential equation alone —
and the corresponding scale factor $a=1/\varepsilon$ is exact de Sitter
expansion, $a(t)=a_0e^{\lambda t}$, with constant Hubble rate $H=\lambda$
and $\ddot a=\lambda^2a>0$: the expansion accelerates for any nonzero
dissipation rate, with no repulsive substance posited anywhere.
``Dark energy'' is accordingly not a substance with negative pressure;
it is this dissipation rate, and matching the observed density gives
$\Lambda \propto \lambda^2$ (specifically $\rho_\Lambda = 3\lambda^2/8\pi G$
once $\lambda$ is identified with the asymptotic Hubble rate) — the
\emph{square} of the dissipation rate, since $\Lambda$ enters through the
acceleration $\ddot a$ rather than through $\lambda$ itself, not the
dissipation rate directly.

Pure exponential dissipation is the sharpest and most falsifiable case,
not the only one available: generalizing $\dot\varepsilon=-\lambda\varepsilon$
to a power law $\dot\varepsilon=-\lambda\varepsilon^q$ for a single
dimensionless exponent $q$ fixes the equation of state exactly,
$w=(2q-5)/3$, invertibly, $q=(3w+5)/2$. The case $q=1$ is exactly
$w=-1$ with no freedom at all — pure exponential dissipation — and it is
the only value of $q$ that gives $w=-1$, not one case among several that
happen to agree. This is the framework's sharpest exposure in this
sector: if $w=-1$ is excluded by observation, pure exponential
dissipation is dead, and the surviving theory must have $q\neq1$, read
off the data by the inversion formula. What the power-law generalization
does \emph{not} provide is a mechanism for $w$ to \emph{evolve}: for any
fixed $q$, $w$ is still a constant, since $q$ is a single dimensionless
exponent rather than a function of time, so DESI's~\cite{DESIDR2}
preference for an
evolving $w_0,w_a$ remains a tension with the shape of the construction
itself and not merely with the value $q=1$ — a constant-but-shifted $w$
is available to this framework in principle, a genuinely time-varying
one is not, which is why the DESI tension of Section~\ref{sec:data} is
correctly read as a tension with A5's constancy requirement rather than
with a single fittable number.

\subsection{Light is scale-blind, but redshift is not}

Reciprocity $s_t\cdot s_r=1$ (Section~\ref{sec:gravity}) holds pointwise,
not only in vacuum, so the null cone defined by $g_{tt}g_{rr}=-1$ is
invariant under any global rescaling $s\mapsto u\cdot s$ permitted by A5.
Since a photon's trajectory is defined purely by the null condition, its
propagation is insensitive to the overall scale — the light cone is
scale-blind by the same fiducial covariance that makes any single scale
unobservable. This gives photon dispersion identically zero at every
energy, a stronger and structurally distinct statement than merely
``dispersion smaller than currently measurable,'' and it is the basis for
the Fermi-LAT comparison of Table~\ref{tab:scoreboard}~\cite{FermiLAT2009}.

This appears to contradict the cosmological picture directly above: if
light cannot see the scale field at all, how can a redshift measurement
ever register cosmic dissipation, which is exactly a change in that
field? And Section~\ref{sec:gravity}'s local-covariance result sharpens
the puzzle rather than resolving it, since it proves a spatially
constant cosmic drift changes \emph{nothing} measured locally. Both
results are theorems, so both are true, and the appearance of
contradiction is a sign that two different things were being called by
the same name. They were. A null vector has no length, but assigning it
an \emph{energy} is not a property of the vector alone — it is the
pairing of the vector with an observer's own direction, and an observer
is made of matter, which (unlike light) does carry the scale. The
\emph{causal structure} light defines is scale-blind, exactly as above;
the \emph{energy} an observer assigns to a light ray is not, because
assigning it requires a scale-carrying observer in the first place.
Redshift is consequently not a property of light at all: it is the
\emph{ratio of two observers' scales}, $(u_1/u_2)^2$, with light
contributing nothing to that ratio beyond carrying the comparison from
one event to another. This simultaneously resolves the apparent tension
with local covariance: a redshift measurement is not a local measurement
of the drift, it compares the scale at emission with the scale at
absorption, and that is a scale \emph{difference}, which A5 declares
physical, while local covariance concerns what a single observer measures
at one point — cosmology observes exactly the quantity the axioms permit
and no more.

\subsection{Questions only this framework can ask}\label{ss:native}

Everything so far has answered questions posed elsewhere: does the
theory give Schwarzschild, does it give a graviton-like sector, does it
give an area law. That is the right first test, and every section above
has run it. It is not the right last one, because a framework built on
different primitives should be able to \emph{ask} things the usual
vocabulary cannot pose at all, and whether it can is where it either
earns its keep or is exposed as a paraphrase of existing physics in new
notation. Seven such questions follow from the primitives already in
use — a scale that is a unit, a pattern that collapses to one vector, a
single global drift, the antisymmetric wedge, a threshold density, the
isotropic/anisotropic dichotomy, and one commutator — each pushed for
what it demands to know rather than for what it was built to answer.

\emph{Separation is not primitive; non-parallelism is.} With no space as
a primitive, ``two places'' cannot be assumed — it must be constructed,
and the wedge is the only construction available (Section~\ref{sec:picture}).
The question this opens, unavailable in a theory that starts from a
manifold: what is the observable that measures \emph{degree} of
separation, if not a distance? The wedge answers it, and being
antisymmetric, the framework's own separation is necessarily
\emph{oriented} — a signed area, not a length, which is not in the
standard vocabulary at all.

\emph{There is one switch, not one per interaction.} The isotropic locus
has no rotational label, no charge, no multiplet, and no curvature
simultaneously, because all four have one source (Section~\ref{sec:branch}).
The framework's own question: the usual account has separate energy
scales at which separate interactions become distinct; here there is one
boundary and then refinement. So — is there any label whose appearance
is \emph{not} a refinement of an already-anisotropic pattern? A
genuinely independent switching-on would falsify the single-source claim,
and it is a question no other framework needs to ask, because no other
framework has only one switch to defend (Section~\ref{ss:secondswitch}
runs this test against fourteen decades of data).

\emph{Content belongs to intervals, never to a single scale.} A window
of zero width resolves nothing; with threshold density $\rho$, a window
of width $\Delta$ contains $\rho\Delta$ structures, vanishing with the
window. ``What exists at energy $E$'' is consequently malformed as a
question in this framework: what exists is a property of a scale
\emph{interval}, never of a point. The complementarity this produces —
sharpness in scale opposed to richness in content — is not Heisenberg
transplanted; the conjugate pair is scale-sharpness against
content-count, and the smallest window in which a single structure is
unambiguous is exactly $1/\rho$.

\emph{There is no global field equation, and A6 is why.} A field
equation can be written only if its source obeys a cyclic conservation
law, and A6 says the universe's global source is not conserved — so
there is no field equation for the universe as a whole, only local ones,
which survive because local covariance survives global dissipation
(Section~\ref{sec:gravity}). The framework's own question: if nothing
global is being solved, what fixes the large-scale history at all? The
only candidate is A5 itself, a symmetry rather than a dynamical law —
the large-scale history is fixed by what is \emph{unobservable}, not by
what evolves, and whether that is sufficient is an open question this
framework poses that general relativity, built around a global field
equation, cannot.

\emph{The quantum-gravity experiment is an ordering test, not a
particle search.} The entire quantum correction to geometry is a
commutator of scale gradients (Section~\ref{sec:quantum}), and a
commutator is an order discrepancy: the native experiment is not
``detect a graviton'' but ``compare the scale along two directions, in
both orders, and look for a discrepancy.'' The magnitude of that
discrepancy is set by the scale step and sits far below anything
currently measurable, stated here rather than hidden; what is
distinctive is that the observable this framework proposes is an
\emph{ordering}, not an amplitude, and no particle needs to exist for
the question to be well posed.

\emph{Defects change parity in pairs, with no coupling constant.} The
confined charge's $\mathbb{Z}/2$ label composes multiplicatively, so the
total parity of any collection cannot change under binding or splitting.
This is not a selection rule imposed on a Lagrangian — there is no
Lagrangian here — it is arithmetic in $\mathbb{Z}/2$: is there any
process that changes total parity? The framework says no, categorically,
which is a cleaner statement than a selection rule derived from an
assumed symmetry, since there is no symmetry being assumed and hence
nothing to break.

\emph{There is no maximum density.} A ratio of scale units is itself a
unit — never zero, never infinite, but unbounded in magnitude — and
matter, in the anisotropic sector, \emph{is} anisotropy of the scale, so
there is no maximum matter density, no Planck-scale density cutoff, and
no scale at which the description must be replaced by a qualitatively
different kind of theory. This is the same statement as the absence of a
minimum length (Section~\ref{sec:gravity}), seen from the matter side
rather than the geometry side, and it is the framework's sharpest
structural disagreement with programs that posit a fundamental cutoff.

None of these seven is a translation of a question asked elsewhere; each
follows from a primitive the standard vocabulary lacks. Their honest
status varies and is stated here rather than left implicit: the
one-switch question and the parity question are testable now in
principle; the interval-content complementarity is already tested
through the resonance-width bound of Section~\ref{sec:flow}; the
no-maximum-density question is testable only through its contrapositive
(a confirmed minimum length would falsify it); the separation, global-history,
and ordering-test questions are conceptual claims with no current
measurement attached. None of the seven is a prediction of a number.

%======================================================================
\section{Confrontation with data}\label{sec:data}
%======================================================================

Table~\ref{tab:scoreboard} is the single scoreboard of this paper.

\begin{table*}
\centering
\small
\begin{tabular}{p{0.24\textwidth}p{0.24\textwidth}p{0.26\textwidth}p{0.18\textwidth}}
\toprule
Quantity & This framework & Measured & Status \\
\midrule
\multicolumn{4}{l}{\emph{Derived from the axiom chain (plus A6$'$); certified to the precision quoted}}\\
Solar light deflection & $1.7515''$ & $1.7509 \pm 0.0002''$ & agrees \\
Mercury perihelion advance & $42.989''$ & $42.98 \pm 0.04''$ & agrees \\
PPN $\gamma$ & $1$ (exact) & $|\gamma-1| \lesssim 2\times10^{-5}$~\cite{Bertotti2003} & agrees \\
\midrule
\multicolumn{4}{l}{\emph{Arithmetic check of the A3 reading, not an independent prediction}}\\
Muon time dilation & $64.43\,\mu\mathrm{s}$ & $64.378 \pm 0.026\,\mu\mathrm{s}$ & agrees ($10^{-3}$) \\
Pound--Rebka & $2.455\times10^{-15}$ & $(2.46\pm0.03)\times10^{-15}$ & agrees \\
\midrule
\multicolumn{4}{l}{\emph{Converted; $\lambda$ itself is measured — the one free unit}}\\
Dark energy density & $5.845\times10^{-27}\,\mathrm{kg/m^3}$ & $5.86\times10^{-27}\,\mathrm{kg/m^3}$~\cite{Planck2018} & agrees \\
\midrule
\multicolumn{4}{l}{\emph{Currently live tests}}\\
$w$ non-evolution (forced by A5) & $w=-1$, no running & DESI: $w_0>-1,\,w_a<0$ & \textbf{tension, $2.8$--$4.2\sigma$} \\
$f_0(500)$ pole $\Gamma/m$ & $<1.157$ & CCL $1.234$~\cite{CCL2006} / GKPY $1.090$~\cite{GKPY2011} & \textbf{open} \\
Photon energy-dependent delay & exactly zero & $E_{\rm QG}\gtrsim 7.6\,E_{\rm Pl}$ & passes, distinct prediction \\
Gravitational wave speed & exactly $c$, no dispersion & $|\Delta v/v|\lesssim10^{-15}$ & passes \\
GW polarization count & exactly $2$, no breathing mode & network polarization tests & passes \\
Dipole radiation & identically $0$ & double pulsar $<1.3\times10^{-4}$ & passes \\
Ringdown echoes & reflectivity $\neq 0$ & GW250114 residual $=$ noise & \textbf{open, under pressure} \\
Ergoregion instability & reflectivity bound uncomputed & GW241011 $\chi_1=0.64$, long-lived & \textbf{open, tighter pressure} \\
Ultralight scalar & \textbf{no such slot} & $10^{-13}$--$3\times10^{-12}\,$eV excluded & compatible \\
Area-law analogue & scale ratio never decreases & consistent with observation & passes (not discriminating) \\
Mass spectrum $\mathrm{CV}^2$ & $=1$ (rank one; mixtures need $\geq1$) & $0.875$ & partial \\
Exactly-conserved unconfined charge & exactly one & electric charge, $B{-}L$(?) & undecided, awaits $0\nu\beta\beta$ \\
Black hole entropy & $\propto \ln R$ & --- & unmeasured (area law also unmeasured) \\
\midrule
\multicolumn{4}{l}{\emph{Excluded by the framework itself — the most informative rows}}\\
Equal-ratio mass tower & $\mathrm{CV}=0$ & $0.935$ & \textbf{excluded} \\
Isotropic-sector-only deflection & $0.8758''$ & $1.7509''$ & \textbf{excluded (factor 2)} \\
Isotropic-sector-only precession & $14.33''$ & $42.98''$ & \textbf{excluded (factor 3)} \\
Color $=$ block charge & $\mathbb{Z}/2$ (two values) & needs three values & \textbf{excluded (structurally)} \\
\bottomrule
\end{tabular}
\caption{The complete numerical scoreboard.}
\label{tab:scoreboard}
\end{table*}

One live tension and one live open question are not fatal — both sectors
are under active measurement — and the last four rows are options the
framework itself rejects; these are more informative than the passing
rows because they demonstrate the framework is capable of being wrong
internally. On timescale: even a seesaw-based lepton-number-violating
model predicts a $0\nu\beta\beta$ half-life far above current exclusion
limits~\cite{KamLANDZen2023}, so this is a genuine but slow-fused test.

\subsection{Equation of state is out of scope}\label{ss:eos-out}

The magnitude of the ringdown reflectivity $\mathcal{R}$
(Section~\ref{ss:echo}) requires an equation of state for the interior
structure replacing a horizon, which varies enormously between candidate
compact objects and which this framework does not attempt to compute —
this is explicitly outside its scope, not a gap it intends to close. What
the framework \emph{does} deliver without an equation of state is the
qualitative echo/ergoregion structure of Section~\ref{ss:echo} and the
comparison between the linear and exponential strength of the two
observational bounds.

%======================================================================
\section{How to falsify this framework}\label{sec:falsify}
%======================================================================

\begin{table}[h]
\centering
\small
\begin{tabular}{p{0.24\linewidth}p{0.24\linewidth}p{0.44\linewidth}}
\toprule
Observation & Kills & Status \\
\midrule
\multicolumn{3}{l}{\textbf{Kills the whole framework}}\\
$w$ evolves with redshift ($5\sigma$, dataset-consistent) & A5, dimensionlessness, all geometry & DESI $2.8$--$4.2\sigma$, but BAO alone self-consistent \\
Energy-dependent photon delay, any order & A3 + light-cone scale-blindness & $E_{\rm QG}\gtrsim7.6E_{\rm Pl}$, winning \\
GW dipole radiation & the conservation theorem, hence Bianchi & double pulsar $<1.3\times10^{-4}$ \\
GW breathing mode & A5 (a breathing mode is a fiducial shift) & tensor favored over scalar/vector/mixed \\
\bottomrule
\end{tabular}
\caption{No adjustable parameter absorbs any of these four: a detection is fatal, not a call for retuning.}
\end{table}

\begin{table}[h]
\centering
\small
\begin{tabular}{p{0.24\linewidth}p{0.28\linewidth}p{0.40\linewidth}}
\toprule
Observation & Kills & Status \\
\midrule
\multicolumn{3}{l}{\textbf{Kills one sector}}\\
A minimum length exists & horizon, singularity and max-density results & same bound as photon dispersion \\
$f_0(500)$ pole $\Gamma/m>1.157$ confirmed & Section~\ref{sec:particles}'s resolvability bound & CCL above, GKPY below \\
$B{-}L$ gauged and exactly conserved & the ``exactly one charge'' result & $0\nu\beta\beta$ not seen, slow fuse \\
A label opens that is not a refinement & Section~\ref{ss:secondswitch}'s single switch & no counterexample in 14 decades \\
Total-parity-violating process & the $\mathbb{Z}/2$ arithmetic & none known \\
\bottomrule
\end{tabular}
\end{table}

\begin{table}[h]
\centering
\small
\begin{tabular}{p{0.28\linewidth}p{0.22\linewidth}p{0.42\linewidth}}
\toprule
Observation & Exposes & Status \\
\midrule
\multicolumn{3}{l}{\textbf{Forces an admission of what cannot be computed}}\\
Rapid spinning compact object decays via ergoregion instability & horizon-absence requires small reflectivity, uncomputed & GW241011 $\chi_1=0.64$, long-lived; tightest current pressure \\
Ringdown echoes seen & same, reversed: predicted to exist, amplitude not predicted & null searches so far \\
BH entropy $\propto$ area, not $\ln R$ & the log-entropy result & entropy itself unmeasured \\
\bottomrule
\end{tabular}
\end{table}

\textbf{A7 cannot be falsified by this table.} It is a commitment about
what counts as an observation; anyone may reject it and obtain a
different $k$, and no experiment adjudicates that choice — which is
precisely why it was promoted to an axiom rather than left an informal
judgment. The one free unit is likewise not falsifiable: it is the scale
in which predictions are expressed, not itself a prediction.

What does \emph{not} constitute falsification: the absence of a Standard
Model particle table (Section~\ref{sec:particles} argues this was never
a goal); the absence of computed hadron masses (the same free unit
again); and, emphatically, agreement with general relativity in the
gravitational-wave sector — Theorem~\ref{thm:allorder} proves the
symmetric sector agrees at every order, so no discrepancy is expected
there in principle.

%======================================================================
\section{Register: retracted, corrected, cited, and assumed}\label{sec:register}
%======================================================================

A theory should carry its own history of correction, not erase it. This
section is a standing answer to the question ``does this rest on the
axioms alone?''

\textbf{Retracted (4).} A geometric mass-tower law, later shown to be an
unproved assumption rather than a derived consequence, and refuted by
lattice glueball data under that assumption; a color identification that
required three values where the projective order parameter structurally
provides only two; a proposed cross-sector consistency test comparing two
quantities that were not, on inspection, the same quantity; a
codimension analysis that described a circle-valued rather than
directional scale.

\textbf{Corrected in scope (6).} In each case the underlying result
survives and the physical interpretation attached to it does not — most
notably, an earlier reading in which the rotational and scaling sectors
were treated as independent channels is true only of the scalar version
of the theory; the directional scale makes the rotation sector
\emph{generated by} the scaling sector (Theorem~\ref{thm:coupling}).

\textbf{A recurring error (5 occurrences).} Five of the corrections above
share one root cause: scalar reasoning applied where the structure is
directional. The diagnostic that catches it: whenever a claim mentions a
count — one axis, one scale, one direction — ask a count of \emph{what};
value and direction are never the same count.

\textbf{Cited, not proved (2).} The classification of real-rank-one real
simple Lie algebras invoked in Section~\ref{sec:chain}, and the relevant
homotopy facts about projective spaces used in the particle-charge
construction of Section~\ref{sec:particles}. If either citation fails,
the weaker statement proved before invoking it still stands; only the
identification built on top does not.

\textbf{Assumed, registered rather than smuggled (3).} A6$'$, the scale
response law, external to the seven axioms; a uniform-density assumption
needed for one dimension-counting step in Section~\ref{sec:chain}; A7
itself, promoted to an axiom rather than left implicit.

\textbf{Structures now witnessed (6 of 6).} Six structures carrying
load-bearing theorems — most importantly the scale-shift structure
carrying Theorem~\ref{thm:hbar} — were, as found by external audit, never
constructed from A1–A7 anywhere in the original development, making the
associated theorems true but conditional. All six are now witnessed by
explicit models: the polynomial-ring construction of
Section~\ref{sec:quantum} witnesses the scale-shift structure itself and
a cosmic-history example; the remaining four are witnessed by separate,
unrelated concrete models — a matrix trace, an adjoint action on a
general ring, a real-valued translation flow, and a constant function —
chosen for whichever structure each was needed for, not by one uniform
construction. This shows the structures are \emph{compatible} with the
axioms, not that the axioms \emph{force} them — the gap between
compatible and forced is real and is not closed by exhibiting a witness.

\textbf{What rests on the axioms alone.} The scale/rotation split, the
codimension-one signature split, the bookkeeping form, the coupling
(Theorem~\ref{thm:coupling}), real rank one (subject to the physical
identification just above), conservation from Bianchi, the full gravity
chain through $1.7515''$ and $42.989''/$century, the label structure of
the anisotropic sector, every statement in Section~\ref{ss:native}
including the one-arrow-of-time identity (Theorem~\ref{thm:onearrow}).
$\hbar$ as an exact scale step is deliberately absent from this list: it
is witnessed, not forced, exactly as stated above.

%======================================================================
\section{Conclusion}\label{sec:concl}
%======================================================================

Taking scale rather than spacetime as primitive, a single chain runs from
axioms to numbers. Scale must be directional, since a scalar version
forbids the Schwarzschild metric outright. Once directional, the
scale–rotation coupling is forced by the meaning of the two words:
scaling is symmetric, rotation is antisymmetric, and the bracket of two
symmetric operators is necessarily antisymmetric. The coupling, together
with the drift being a single functional, gives real rank one, hence the
Lorentz family, hence a fixed dimension with three spatial directions set
by A7 — a commitment about what counts as an observation, not a free
parameter. The metric connection and the gauge connection turn out to be
the same commutator, read under two different questions. The
isotropic/anisotropic branch determines whether labels exist at all,
itself a scale-dependent fact, giving this framework's account of why
interactions separate as energy falls. Curvature is scale deformation,
the renormalization group is the same drift read as a running coupling,
$\hbar$ is a step along the scale axis, and the entire quantum correction
to geometry is purely antisymmetric. Dark matter is the isotropic
trajectory of the same field that produces ordinary gravity, and dark
energy is that field's own dissipation rate. The gravity sector closes
quantitatively at $1.7515''$ and $42.989''/$century with nothing assumed
in between.

These are not metaphors; every claim stated as a theorem in this paper
has an independent, complete proof, verified end to end by machine — 644
checked theorems, no admitted gaps, and no hidden axioms beyond those of
standard mathematics. But the statement that this framework reproduces
known physics cannot currently be made without qualification: the
scoreboard carries one live tension, DESI's $2.8$–$4.2\sigma$ preference
against $w$ non-evolution, which is a tension with A5 itself and not with
a fittable parameter, since no version of this framework admits an
evolving $w$; and one live open question, the side of the resolvability
bound on which the true $f_0(500)$ pole lies. Nor does the framework
explain the Standard Model — Section~\ref{sec:particles} argues that was
never the goal — and Section~\ref{sec:register} records four retracted
claims, six claims corrected in scope, two citations of results not
reproved here, and three assumptions registered rather than hidden, along
with one recurring scalar-for-directional error caught five times.

What a fully verified argument can establish is exactly one thing: that
these conclusions follow from these axioms, with no gap in between.
Whether nature adopts the axioms is a separate question that this paper
does not, and should not be read to, answer — which is exactly why the
scoreboard and the register are given in full.

\begin{acknowledgments}
The complete formal development and its supporting tooling are available
in the project repository, along with a companion technical report.
\end{acknowledgments}

\onecolumngrid
\appendix
\section{Note on the Supplemental Material}
The \SM\ gives, in the same order as the sections above, a complete
mathematical restatement of every one of the 644 machine-verified
theorems underlying this paper — full formal statements and full proofs,
organized to mirror the structure of this paper section by section, with
each item individually identifiable for cross-reference. This paper does
not depend on the \SM\ to be read; the \SM\ exists so that every claim
made informally above can be checked in complete formal detail.

\bibliography{refs}

\end{document}
